
% Authoritative LaTeX manuscript; Markdown is a generated compatibility projection.

The dashboard asks whether autonomous work is producing durable, trustworthy progress. Read the outcomes first; use diagnostics to explain why they changed.

Track five primary readings: durable throughput, defect escape, human-attention burden, representation health, and engineering-capital return. The first three measure outcomes; the last two ask whether the environment can sustain them. Churn, model coverage, support ratio, control count, grammar coverage, and navigation cost are diagnostics, not objectives.

{\footnotesize
\begin{longtable}[]{@{}
  >{\raggedright\arraybackslash}p{(\linewidth - 6\tabcolsep) * \real{0.2500}}
  >{\raggedright\arraybackslash}p{(\linewidth - 6\tabcolsep) * \real{0.2500}}
  >{\raggedright\arraybackslash}p{(\linewidth - 6\tabcolsep) * \real{0.2500}}
  >{\raggedright\arraybackslash}p{(\linewidth - 6\tabcolsep) * \real{0.2500}}@{}}
\caption[The Operator\textquotesingle s Dashboard --- five primary readings]{The Operator\textquotesingle s Dashboard. Five primary readings, the question each answers, the desired direction of change, and useful diagnostics.}\label{operators-dashboard}\tabularnewline
\toprule\noalign{}
\begin{minipage}[b]{\linewidth}\raggedright
Reading
\end{minipage} & \begin{minipage}[b]{\linewidth}\raggedright
What to ask
\end{minipage} & \begin{minipage}[b]{\linewidth}\raggedright
Healthy direction
\end{minipage} & \begin{minipage}[b]{\linewidth}\raggedright
Useful diagnostics
\end{minipage} \\
\midrule\noalign{}
\endfirsthead
\toprule\noalign{}
\begin{minipage}[b]{\linewidth}\raggedright
Reading
\end{minipage} & \begin{minipage}[b]{\linewidth}\raggedright
What to ask
\end{minipage} & \begin{minipage}[b]{\linewidth}\raggedright
Healthy direction
\end{minipage} & \begin{minipage}[b]{\linewidth}\raggedright
Useful diagnostics
\end{minipage} \\
\midrule\noalign{}
\endhead
\bottomrule\noalign{}
\endlastfoot
\textbf{Durable throughput} & How much accepted work survives without reopening, rollback, or repeated repair? & More accepted capability without proportional growth in intervention & closure rate; reopen rate; change cadence; churn \\
\textbf{Defect escape} & What incorrect or policy-violating work crosses the boundary that was supposed to catch it? & Falls for governed obligations & validator findings; escaped defects; rollback/incident rate \\
\textbf{Human-attention burden} & How much repeated reconstruction, review, or adjudication does each unit of durable work require? & Falls where knowledge or judgment can be made durable; remains where human judgment is deliberate & interventions per landed change; review time; recurring decisions \\
\textbf{Representation health} & Can the representations being relied upon still answer the questions they claim to answer? & Claimed correspondence holds; stale or uncovered surfaces remain explicit & drift checks; traceability; freshness; coverage/relevance \\
\textbf{Engineering-capital return} & Does accumulated structure make later work more capable or cheaper than the structure costs to carry? & Reuse and inherited capability rise while carrying cost remains justified & mechanism reuse; recurring-class disappearance; maintenance burden; retirement candidates \\
\end{longtable}
}

DocAble\textquotesingle s support ratio, control count, Missing-Model drain, navigation pilot, and churn curves are observations from one build, not operating targets. Appendix H records the underlying measurements and their limitations.

Treat uncalibrated visual indicators as status, not quantitative measurements.
